\section{A range-to-root transfer principle}
\label{sec:abstract-response}

The construction of a selected complete-history root has two logically
different parts.  A range argument must first construct the history graph and
the two selected traces, including their parameter derivatives.  Once that
has been done, a scalar argument turns two Melnikov pairings and their
remainders into a root response.  We isolate the latter implication here so
that the model-specific calculation below has a reusable conclusion.  In
particular, the following theorem does not construct an invariant graph or a
trace inverse: their existence and regularity are explicit hypotheses.

Let \(I\subset\R\) be a compact interval.  For each \(N\ge2\), let
\(\mathcal H_N\) be a history space, let
\(\nu_{0,N}\in I\), and assume
\begin{equation}
 \inf_{N\ge2}\operatorname{dist}(\nu_{0,N},\partial I)=d_*>0.
 \label{eq:abstract-center-buffer}
\end{equation}
The notation below permits the graph, the trace spaces, and their norms to
depend on \(N\) and \(\delta\); only the constants displayed as common are
required to be independent of \(N\).

\begin{theorem}[Uniform range-to-root and Melnikov transfer]
\label{thm:abstract-root-response}
Fix \(\bar\delta,\bar\zeta>0\).  Suppose that the following data have been
verified for every \(N\ge2\), \(0<\delta<\bar\delta\), \(\nu\in I\), and
\(|\zeta|<\bar\zeta\).

\smallskip
\noindent
\emph{Complete-history range.}
An already solved range problem supplies a Banach chart \(\mathcal U_N\),
a phase section \(\Sigma_N\subset\mathcal U_N\), an exact
complete-history lift with invariant range
\[
 \mathcal I_{N,\delta,\nu,\zeta}:\mathcal U_N\longrightarrow\mathcal H_N,
\]
and selected one-sided traces
\(z^a_{N,\delta,\nu,\zeta}\) and
\(z^r_{N,\delta,\nu,\zeta}\).  Their central values lie on the same phase
section.  Present-state evaluation \(\operatorname{ev}_{0,N}\) is a left
inverse of \(\mathcal I_{N,\delta,\nu,\zeta}\), so the lift is injective.
A scalar matcher
\(\mathfrak m_{N,\delta,\nu,\zeta}:\Sigma_N\times\Sigma_N\to\R\)
defines
\begin{equation}
 \mathscr G_N(\delta,\nu,\zeta)
 =\mathfrak m_{N,\delta,\nu,\zeta}
   \bigl(z^a_{N,\delta,\nu,\zeta}(0),
         z^r_{N,\delta,\nu,\zeta}(0)\bigr).
 \label{eq:abstract-gap-definition}
\end{equation}
The matcher has the zero-fiber property
\begin{equation}
 \mathscr G_N(\delta,\nu,\zeta)=0
 \quad\Longleftrightarrow\quad
 \mathcal I_{N,\delta,\nu,\zeta}
       \bigl(z^a_{N,\delta,\nu,\zeta}(0)\bigr)
 =\mathcal I_{N,\delta,\nu,\zeta}
       \bigl(z^r_{N,\delta,\nu,\zeta}(0)\bigr).
 \label{eq:abstract-zero-fiber}
\end{equation}
Thus a zero is equality of the selected histories, not merely equality of a
projected observable.  The lift, traces, and matcher possess the parameter
jets needed to make \(\mathscr G_N\) jointly \(C^2\) in \((\nu,\zeta)\) on
an open neighborhood of the closed parameter box.  This includes the pure
\(\nu\nu\) and mixed \(\nu\zeta\) derivatives.

\smallskip
\noindent
\emph{Melnikov data and tails.}
There are Banach spaces \(V_N\), measurable functions
\(\psi_N:\R\to V_N^*\) and
\(f_{\nu,N},f_{\zeta,N}:\R\to V_N\), and numbers
\(R_\delta\to\infty\) as \(\delta\to0\), such that the absolutely
convergent pairings
\begin{align}
 a_N&=\int_{\R}\langle\psi_N(s),f_{\nu,N}(s)\rangle\,\dd s,
 &
 b_N&=\int_{\R}\langle\psi_N(s),f_{\zeta,N}(s)\rangle\,\dd s
 \label{eq:abstract-melnikov-coefficients}
\end{align}
represent the derivatives of the same normalized matcher as in
\cref{eq:abstract-gap-definition}.  Thus any phase, endpoint, or moving-history
contribution must already be included in these profiles or proved to vanish.
Assume that, with constants independent of \(N\),
\begin{equation}
 0<a_*\le |a_N|\le a^*,\qquad |b_N|\le b^*.
 \label{eq:abstract-coefficient-bounds}
\end{equation}
Put
\begin{equation}
 a_{N,\delta}=\int_{-R_\delta}^{R_\delta}
     \langle\psi_N,f_{\nu,N}\rangle\,\dd s,
 \qquad
 b_{N,\delta}=\int_{-R_\delta}^{R_\delta}
     \langle\psi_N,f_{\zeta,N}\rangle\,\dd s.
 \label{eq:abstract-truncated-coefficients}
\end{equation}
Assume the common tail estimate
\begin{equation}
 \int_{|s|>R_\delta}
 \left(
  |\langle\psi_N(s),f_{\nu,N}(s)\rangle|
  +|\langle\psi_N(s),f_{\zeta,N}(s)\rangle|
 \right)\dd s
 \le C_{\rm tail}\delta.
 \label{eq:abstract-tail-bound}
\end{equation}

\smallskip
\noindent
\emph{Range-to-gap remainders.}
The output of the graph and trace estimates obeys, throughout the same
parameter box,
\begin{align}
 \left|\frac{\mathscr G_N(\delta,\nu,\zeta)}{\delta}
       -a_{N,\delta}(\nu-\nu_{0,N})\right|
 &\le C_{\rm rg}\delta,
 \label{eq:abstract-gap-level}\\
 |\partial_\nu\mathscr G_N-\delta a_{N,\delta}|
 &\le C_{\rm rg}\delta^2,
 \label{eq:abstract-gap-slope}\\
 |\partial_\zeta\mathscr G_N-\delta^2b_{N,\delta}|
 &\le C_{\rm rg}(\delta^3+\delta^2|\zeta|),
 \label{eq:abstract-gap-response}\\
 |\partial_{\nu\nu}\mathscr G_N|
 +|\partial_{\nu\zeta}\mathscr G_N|
 +|\partial_{\zeta\zeta}\mathscr G_N|
 &\le C_{\rm rg}\delta^2.
 \label{eq:abstract-gap-second-jets}
\end{align}

Then there are \(\delta_0\in(0,\bar\delta)\),
\(\zeta_0\in(0,\bar\zeta)\), \(c_0\in(0,d_*)\), and \(C>0\), all
independent of \(N\), such that for every \(N\ge2\),
\(0<\delta<\delta_0\), and \(|\zeta|<\zeta_0\), the equation
\(\mathscr G_N(\delta,\nu,\zeta)=0\) has exactly one solution in
\(|\nu-\nu_{0,N}|<c_0\).  This solution is a selected complete-history
root by \cref{eq:abstract-zero-fiber}, is \(C^2\) in \(\zeta\), and
satisfies
\begin{align}
 \partial_\zeta\nu_{c,N}(\delta,\zeta)
 &=-\delta\frac{b_N}{a_N}
   +O(\delta^2+\delta|\zeta|),
 \label{eq:abstract-root-derivative}\\
 |\partial_{\zeta\zeta}\nu_{c,N}(\delta,\zeta)|
 &\le C\delta,
 \label{eq:abstract-root-second-derivative}\\
 \left|\nu_{c,N}(\delta,\zeta)-\nu_{c,N}(\delta,0)
       +\delta\frac{b_N}{a_N}\zeta\right|
 &\le C(\delta^2|\zeta|+\delta\zeta^2).
 \label{eq:abstract-root-increment}
\end{align}
The constant implicit in \cref{eq:abstract-root-derivative} is common in
\(N\).  More generally, if a physical parameter is scaled as
\(\mu=\mu_f(\delta)+\delta^q\nu\), where \(q>0\) and \(\mu_f\) is
independent of \(\zeta\), then
\begin{equation}
 \left|\mu_{c,N}(\delta,\zeta)-\mu_{c,N}(\delta,0)
       +\frac{b_N}{a_N}\delta^{q+1}\zeta\right|
 \le C\{\delta^{q+2}|\zeta|+\delta^{q+1}\zeta^2\}.
 \label{eq:abstract-physical-increment}
\end{equation}
In particular, \(q=2\) gives a leading \(\delta^3\zeta\) response.
\end{theorem}

\begin{proof}
The tail bound gives
\begin{equation}
 |a_{N,\delta}-a_N|+|b_{N,\delta}-b_N|
 \le C_{\rm tail}\delta.                            \label{eq:abstract-tail-consequence}
\end{equation}
Consequently, after decreasing \(\delta_0\),
\begin{align}
 \left|\frac{\mathscr G_N}{\delta}
       -a_N(\nu-\nu_{0,N})\right|&\le C_0\delta,
 \label{eq:abstract-full-gap-level}\\
 |\partial_\nu\mathscr G_N-\delta a_N|&\le C_1\delta^2,
 \label{eq:abstract-full-gap-slope}\\
 |\partial_\zeta\mathscr G_N-\delta^2b_N|
 &\le C_2(\delta^3+\delta^2|\zeta|),
 \label{eq:abstract-full-gap-response}
\end{align}
uniformly in the parameter box.  Here the bounded diameter of \(I\) absorbs
the factor \(|\nu-\nu_{0,N}|\) when
\cref{eq:abstract-tail-consequence} is inserted into
\cref{eq:abstract-gap-level}.

Set \(\zeta_0=\bar\zeta/2\) and choose, for example,
\(c_0=d_*/2\).  At
\(\nu=\nu_{0,N}\pm c_0\),
\cref{eq:abstract-full-gap-level} has the two opposite leading values
\(\pm a_Nc_0\).  They retain opposite signs once
\(C_0\delta_0<a_*c_0/2\), so continuity gives a zero in the local
window.  If also \(C_1\delta_0<a_*/2\), then
\begin{equation}
 |\partial_\nu\mathscr G_N|\ge\frac12a_*\delta
 \label{eq:abstract-uniform-slope}
\end{equation}
and its sign is the sign of \(a_N\) throughout that window.  The zero is
therefore unique.  The ordinary implicit-function theorem, applied at each
\(\zeta\), and uniqueness patch these local branches into one \(C^2\)
function on the common \(\zeta\)-interval.

Along the root,
\begin{equation}
 \partial_\zeta\nu_{c,N}
 =-\frac{\partial_\zeta\mathscr G_N}
          {\partial_\nu\mathscr G_N}.                \label{eq:abstract-ift-ratio}
\end{equation}
Write the numerator and denominator as
\(\delta^2b_N+e_\zeta\) and \(\delta a_N+e_\nu\).  By
\cref{eq:abstract-full-gap-slope,eq:abstract-full-gap-response},
\(|e_\nu|\le C_1\delta^2\) and
\(|e_\zeta|\le C_2(\delta^3+\delta^2|\zeta|)\).  Using
\cref{eq:abstract-uniform-slope,eq:abstract-coefficient-bounds} in
\cref{eq:abstract-ift-ratio} gives
\begin{align*}
 \left|\partial_\zeta\nu_{c,N}
       +\delta\frac{b_N}{a_N}\right|
 &\le \frac{|e_\zeta|}{|\delta a_N+e_\nu|}
  +\frac{\delta^2|b_N|\,|e_\nu|}
   {|\delta a_N|\,|\delta a_N+e_\nu|} \\
 &\le C(\delta^2+\delta|\zeta|),
\end{align*}
which is \cref{eq:abstract-root-derivative}.

Differentiating
\(\mathscr G_N(\delta,\nu_{c,N}(\zeta),\zeta)=0\) once more gives
\begin{equation*}
 \partial_\nu\mathscr G_N\,\partial_{\zeta\zeta}\nu_{c,N}
 =-\left\{
   \partial_{\nu\nu}\mathscr G_N(\partial_\zeta\nu_{c,N})^2
  +2\partial_{\nu\zeta}\mathscr G_N\partial_\zeta\nu_{c,N}
  +\partial_{\zeta\zeta}\mathscr G_N\right\}.
\end{equation*}
The first derivative is \(O(\delta)\); hence
\cref{eq:abstract-gap-second-jets,eq:abstract-uniform-slope} yield
\cref{eq:abstract-root-second-derivative}.  Integrating
\cref{eq:abstract-root-derivative} from \(0\) to \(\zeta\) proves
\cref{eq:abstract-root-increment}.  Multiplication by \(\delta^q\) proves
\cref{eq:abstract-physical-increment}.
\end{proof}

\paragraph{How to apply the theorem.}
For a new RFDE or other infinite-dimensional matching problem, one must
verify five items: (i) construct one exact history lift and the two selected
traces on it; (ii) prove the zero-fiber property of the scalar matcher;
(iii) justify the joint \(C^2\) graph--trace parameter jets, including the
pure and mixed derivatives in \cref{eq:abstract-gap-second-jets};
(iv) identify the adjoint pairings \(a_N,b_N\), including all phase and
endpoint terms, prove
\(\inf_N|a_N|>0\), and estimate their finite-section tails; and (v) prove
the four range-to-gap bounds
\cref{eq:abstract-gap-level,eq:abstract-gap-slope,%
eq:abstract-gap-response,eq:abstract-gap-second-jets}.  A nonzero response
additionally requires a model-specific bound on \(|b_N/a_N|\); it is not a
consequence of the transfer theorem.

\paragraph{Structural and model-specific inputs.}
The exact history lift, selected trace range, zero-fiber implication,
parameter jets, transverse inverse estimates used to obtain them, and the
tail and remainder bounds are the structural inputs to
\cref{thm:abstract-root-response}.  The theorem itself is indifferent to
whether they arise from an RFDE graph transform, a Lin problem, or another
range construction.  In the shared-resource network, by contrast, the
identities
\[
 X_0(s)-X_0(s-\theta_k)=-\frac{\theta_k}{2\alpha},
 \qquad
 \psi(s)=e^{-s^2/2}(s,1)^T
\]
and the stationary-row cancellation are specific to the affine singular
orbit and the chosen delay layers.  The curvature--resolvent return then
gives
\begin{equation}
 a_N=\sqrt{2\pi},\qquad
 b_N=\sqrt{2\pi}\frac{r_N}{\alpha},\qquad
 -\frac{b_N}{a_N}=-\frac{r_N}{\alpha}=\mathscr C_N.
 \label{eq:abstract-current-dictionary}
\end{equation}
The graph and trace hypotheses are supplied in
\cref{lem:shared-graph-trace}; the exact finite-section passage is
\cref{prop:finite-gap-bridge}.  Together with
\cref{eq:shared-leading-gap,eq:shared-melnikov,%
eq:shared-gaussian-tail,eq:shared-gap-expansion,eq:shared-gap-nu,%
eq:shared-gap-zeta,eq:shared-gap-zz,eq:shared-gap-mixed}
verify the pairing and remainder inputs.  In another model, these affine,
Gaussian, and curvature identities must be recomputed; the abstract theorem
does not propagate them automatically.
